\headshort{Онол, технологи}
\chapter{Онол, технологийн судалгаа}
\section{Хүний байрлал тодорхойлох технологи}

Хүний байрлал тодорхойлох (Human Pose Estimation, HPE) нь зураг эсвэл видео бичлэгээс хүний биеийн гол цэгүүдийн (keypoint) байрлалыг тодорхойлох компьютер харааны даалгавар юм. Гол цэг гэдэг нь хамар, нүд, мөр, тохой, бугуй, гуя, өвдөг, шагай зэрэг биеийн анатомийн тодорхой цэгүүдийг хэлнэ. Эдгээр цэгүүдийг хооронд нь шугамаар холбоход хүний биеийн араг ясны (skeleton) энгийн загвар үүсдэг бөгөөд энэ загвар нь хүний байрлал, хөдөлгөөнийг тоон хэлбэрээр илэрхийлэх боломжийг олгодог.

HPE загварын гаралт нь гол цэг бүрийн хувьд координат болон итгэлцлийн оноо (confidence score) байдаг. Итгэлцлийн оноо нь тухайн цэгийг хэр найдвартай илрүүлснийг 0-1 хооронд илэрхийлэх бөгөөд биеийн хэсэг кадраас гарсан, бүдэг харагдсан үед бага утгатай байдаг.

HPE аргуудыг дараах шинжээр ангилдаг. Үүнд:

\begin{itemize}
\item
  \textbf{Гаралтын хэмжээсээр:} 2D аргууд нь зурган дээрх (x, y) пикселийн координатыг, 3D аргууд нь гүний мэдээлэл бүхий (x, y, z) координатыг тодорхойлно. Энгийн нэг камераар 3D байрлалыг тодорхойлох нь гүний мэдээлэл дутмаг тул илүү төвөгтэй бөгөөд алдаа ихтэй байдаг.
\item
  \textbf{Хүний тоогоор:} Нэг хүний (single-person) аргууд зурагт ганц хүн байна гэж үздэг бол олон хүний (multi-person) аргууд нэг кадрт байгаа бүх хүнийг илрүүлдэг. Энэхүү ажилд хоёр болон түүнээс дээш хүнийг нэгэн зэрэг боловсруулах шаардлагатай тул олон хүний арга шаардлагатай.
\item
  \textbf{Олон хүнийг боловсруулах зарчмаар:} Дээрээс доош (top-down) аргууд эхлээд хүн бүрийг хүрээгээр (bounding box) илрүүлж, дараа нь хүрээ тус бүрт гол цэгүүдийг тодорхойлдог. Доороос дээш (bottom-up) аргууд эхлээд зураг дээрх бүх гол цэгийг илрүүлж, дараа нь тэдгээрийг хүн тус бүрээр бүлэглэдэг. Дээрээс доош аргууд ихэвчлэн илүү нарийвчлалтай боловч тооцооллын хугацаа нь хүний тоотой пропорциональ өсдөг.
\item
  \textbf{Гол цэгийг тодорхойлох аргаар:} Дулааны зурагт (heatmap) суурилсан аргууд гол цэг бүрт магадлалын зураг үүсгэж, хамгийн их утгатай цэгийг сонгодог. Регрессэд суурилсан аргууд координатыг шууд тооцоолдог бөгөөд ихэвчлэн хурдан боловч нарийвчлал бага байж болно.
\end{itemize}

\section{Гол цэгийн формат ба тоо}

HPE загварууд нь сургалтад ашигласан өгөгдлийн сангаас (COCO \cite{lin2014}, MPII \cite{andriluka2014}, COCO-WholeBody \cite{jin2020}) хамааран өөр өөр тооны гол цэгийг илрүүлдэг. Биеийн үндсэн хэсгийг л илрүүлдэг загварууд 14-17 цэг гаргадаг бол гар, хуруу, нүүрний хэлбэрийг нарийвчлан илрүүлдэг загварууд хэдэн зуун цэг гаргадаг.

\begin{table}[H]
\centering
\caption{Гол цэгийн формат ба тоо}
\label{tab:formats}
\small
\begin{tabular}{|>{\raggedright\arraybackslash}p{0.234\linewidth}|>{\raggedright\arraybackslash}p{0.260\linewidth}|>{\raggedright\arraybackslash}p{0.372\linewidth}|}
\hline
\textbf{Формат / загвар} & \textbf{Гол цэгийн тоо} & \textbf{Хамрах хүрээ} \\ \hline
LSP & 14 & Зөвхөн бие \\ \hline
MPII & 16 & Зөвхөн бие \\ \hline
COCO (YOLO-Pose, HRNet, RTMPose) & 17 & Бие, нүүрний 5 цэг; хөлийн ул, хуруугүй \\ \hline
OpenPose BODY\_25 & 25 & Бие ба хөлийн ул \\ \hline
MediaPipe Pose (BlazePose) & 33 & Бие, хөлийн ул, гар болон нүүрний зарим цэг \\ \hline
COCO-WholeBody & 133 & Бие, хөлийн ул, хуруу, нүүр \\ \hline
Halpe & 136 & Бие, хөлийн ул, хуруу, нүүр \\ \hline
MediaPipe Holistic & 543 & Бие (33), нүүрний тор (468), гар тус бүр (21) \\ \hline
\end{tabular}
\end{table}

Гол цэгийн тоо их байх нь үргэлж давуу тал биш юм. Цэгийн тоо нэмэгдэхийн хэрээр загварын тооцооллын хугацаа уртасдаг бөгөөд хуруу, нүүрний жижиг цэгүүд нь бага нягтаршилтай бичлэгт алдаа ихтэй илэрдэг. Иймд загвар сонгохдоо тухайн даалгаварт хэдэн гол цэг шаардлагатайг эхлээд тодорхойлох нь зүйтэй. Биеийн ерөнхий хөдөлгөөний синхрончлолыг үнэлэхэд мөр, тохой, бугуй, гуя, өвдөг, шагайн цэгүүд гол үүрэгтэй тул 17-33 цэг хангалттай байдаг. Хурууны хөдөлгөөн эсвэл нүүрний хувирлыг үнэлэх шаардлагатай тохиолдолд л бүтэн биеийн (whole-body) форматыг сонгоно.

\headshort{HPE загварууд}
\section{Хүний байрлал тодорхойлох загварууд}

\textbf{OpenPose} \cite{cao2017} нь доороос дээш зарчмаар ажилладаг анхны өргөн хэрэглэгдсэн олон хүний HPE систем юм. Энэ загвар нь гол цэгүүдийн дулааны зургаас гадна хэсгийн холбоосын талбар (Part Affinity Fields, PAF) хэмээх вектор талбар тооцоолж, түүгээр гол цэгүүдийг хүн тус бүрээр холбодог. Хүний тоо нэмэгдэхэд тооцооллын хугацаа бараг өөрчлөгддөггүй давуу талтай боловч орчин үеийн загваруудаас удаан, их тооцооллын нөөц шаарддаг.

\textbf{HRNet} (High-Resolution Network) \cite{sun2019} нь дээрээс доош зарчмаар ажилладаг бөгөөд зургийн өндөр нягтаршилтай дүрслэлийг сүлжээний туршид хадгалж, өөр өөр нягтаршилтай салбаруудын мэдээллийг тасралтгүй нэгтгэдэг. Үүний үр дүнд гол цэгүүдийг өндөр нарийвчлалтай тодорхойлдог боловч тооцоолол ихтэй тул бодит хугацаанд ажиллуулахад хүчирхэг GPU шаардлагатай.

\textbf{MediaPipe Pose (BlazePose)} \cite{bazarevsky2020} нь Google-ийн боловсруулсан, гар утас болон CPU дээр бодит хугацаанд ажиллахаар оновчлогдсон хөнгөн загвар юм. 33 гол цэг гаргадаг бөгөөд цэг бүрийн хувьд харьцангуй гүний (z) утга болон харагдах магадлалыг өгдөг. Анхандаа нэг хүнд зориулагдсан байсан бөгөөд шинэ хувилбарын Pose Landmarker нь хэд хэдэн хүнийг илрүүлэх тохиргоотой.

\textbf{YOLO-Pose} \cite{maji2022} нь объект илрүүлэх YOLO загварын бүтцэд гол цэгийн тооцооллыг нэмж, хүнийг илрүүлэх болон гол цэгийг тодорхойлохыг нэг алхамд гүйцэтгэдэг. Ultralytics-ийн \cite{ultralytics} YOLOv8-Pose, YOLO11-Pose хувилбарууд COCO форматын 17 цэг гаргадаг бөгөөд хурд ба нарийвчлалын сайн тэнцвэртэй, хэрэглэхэд хялбар байдаг.

\textbf{RTMPose} \cite{jiang2023} нь MMPose сангийн бодит хугацааны дээрээс доош загвар бөгөөд COCO болон COCO-WholeBody форматыг дэмждэг. Хурд, нарийвчлалын хувьд орчин үеийн өрсөлдөхүйц үр дүн үзүүлдэг.

\textbf{ViTPose} \cite{xu2022} нь зургийн трансформер (Vision Transformer) архитектурт суурилсан дээрээс доош загвар бөгөөд энгийн бүтэцтэй хэрнээ гол цэгийг өндөр нарийвчлалтай тодорхойлдог. Трансформерийн тооцоолол их тул бусад загвараас харьцангуй удаан ажилладаг.

\begin{table}[H]
\centering
\caption{HPE загваруудын харьцуулалт}
\label{tab:models}
\small
\begin{tabular}{|>{\raggedright\arraybackslash}p{0.116\linewidth}|>{\raggedright\arraybackslash}p{0.132\linewidth}|>{\raggedright\arraybackslash}p{0.124\linewidth}|>{\raggedright\arraybackslash}p{0.155\linewidth}|>{\raggedright\arraybackslash}p{0.248\linewidth}|}
\hline
\textbf{Загвар} & \textbf{Зарчим} & \textbf{Гол цэг} & \textbf{Олон хүн} & \textbf{Онцлог} \\ \hline
OpenPose & Доороос дээш & 25 (+ гар, нүүр 135) & Тийм & Удаан, их нөөц шаарддаг \\ \hline
HRNet & Дээрээс доош & 17 / 133 & Тийм (илрүүлэгчтэй) & Өндөр нарийвчлал, GPU шаарддаг \\ \hline
MediaPipe Pose & Дээрээс доош & 33 & Хязгаарлагдмал & CPU дээр хурдан, хөнгөн \\ \hline
YOLO-Pose & Нэг алхамт & 17 & Тийм & Хурд, нарийвчлалын тэнцвэр сайн \\ \hline
RTMPose & Дээрээс доош & 17 / 133 & Тийм (илрүүлэгчтэй) & Хурдан, өндөр нарийвчлал \\ \hline
ViTPose & Дээрээс доош & 17 / 133 & Тийм (илрүүлэгчтэй) & Трансформерт суурилсан, өндөр нарийвчлал, удаан \\ \hline
\end{tabular}
\end{table}

\headshort{Үений өнцгийн дүрслэл}
\section{Үений өнцөгт суурилсан байрлалын дүрслэл}

HPE загварын гаргасан координатууд пикселийн нэгжээр илэрхийлэгддэг тул хүний зураг дээрх байрлал, хэмжээнээс шууд хамаардаг. Хоёр хүн ижил хөдөлгөөн хийж байгаа ч нэг нь камерт ойр, нөгөө нь хол зогсвол ойр хүний араг яс том, хол хүнийх жижиг харагдах бөгөөд кадрын өөр хэсэгт зогссон тул координатууд нь шилжилттэй байна. Координатыг шууд харьцуулах тохиолдолд эдгээр нөлөөг шилжилт ба масштабын нормчлол эсвэл Прокрустын шинжилгээгээр арилгах шаардлагатай болдог. Харин энэхүү ажилд байрлалыг үений өнцгөөр дүрсэлдэг бөгөөд үений өнцөг нь хүний кадр дахь байрлал болон зураг дээрх хэмжээнээс угаасаа хамааралгүй тул тусдаа нормчлол шаардлагагүй.

\textbf{Үений өнцөг.} Байрлалыг координатаар биш, үений өнцгөөр илэрхийлэх нь масштаб болон шилжилтээс угаасаа хамааралгүй арга юм. Жишээлбэл, тохойн өнцгийг мөр a, тохой b, бугуй c цэгүүдээр дараах байдлаар тооцоолно:

\begin{equation}
\theta = \arccos \left( \frac{(a - b) \cdot (c - b)}{\|a - b\| \, \|c - b\|} \right)
\end{equation}

Үений өнцгийн хязгаарлалт нь 2D зураг дээр тооцоологддог тул камерын харах чиглэлээс хамаардаг явдал юм. Нэг хүн камерт нүүрээрээ, нөгөө нь хажуугаараа харж зогсвол ижил хөдөлгөөн өөр өнцөгтэй харагдана. Иймд бичлэгт хүмүүс камерт ойролцоо ижил чиглэлээр харсан байх нөхцөл тавина (4.1). Мөн камераас хол байгаа хүний гол цэгүүд илүү чичирдэг тул өнцөг шуугиантай болох бөгөөд үүнийг урьдчилсан боловсруулалтын тэгшитгэх алхмаар бууруулна (3.4).

\textbf{Толин тусгал хөдөлгөөн.} Бүжгийн найруулгаас хамааран хоёр хүн зарим үед толин тусгал хэлбэрээр хөдөлдөг. Жишээлбэл, нэг бүжигчин зүүн гараа, нөгөө нь баруун гараа яг зэрэг өргөдөг. Ердийн харьцуулалт нь нэг хүний зүүн гарыг нөгөөгийн зүүн гартай харьцуулдаг тул ийм хөдөлгөөн бүрэн синхрон байсан ч бага оноо авна. Толин тусгал харьцуулалтад нэг хүний зүүн, баруун талын гол цэгүүдийн индексийг сольж, жишээлбэл А хүний зүүн тохойн өнцгийг Б хүний баруун тохойн өнцөгтэй харьцуулна. Үений өнцөг нь толин тусгалаар өөрчлөгддөггүй тул координатыг эргүүлэх шаардлагагүй, зөвхөн индекс солиход хангалттай.

\headshort{Синхрончлолын хэмжүүрүүд}
\section{Хөдөлгөөний ижил төстэй байдал ба синхрончлолын хэмжүүрүүд}

Хөдөлгөөний синхрончлолыг үнэлэх хэмжүүрүүдийг тухайн нэг агшны байрлалын ижил төстэй байдлыг хэмжих (орон зайн) ба цаг хугацааны дагуух хөдөлгөөний уялдааг хэмжих (хугацааны) гэсэн хоёр төрөлд хуваана.

\textbf{Үений өнцгийн зөрүү.} Хоёр хүний харгалзах үений өнцгүүдийн абсолют зөрүүний дунджаар байрлалын ялгааг үнэлнэ. Энэ хэмжүүр нь шууд тайлбарлах боломжтой (жишээлбэл, «тохойн өнцөг дунджаар 12° зөрүүтэй») давуу талтай.

\textbf{Пирсоны корреляцийн коэффициент.} Хоёр хүний тодорхой нэг үзүүлэлтийн (жишээлбэл, тохойн өнцгийн) хугацааны цуваа хэр зэрэг хамт өөрчлөгдөж байгааг -1-ээс 1 хүртэлх утгаар илэрхийлнэ:

\begin{equation}
r = \frac{\sum_{t} (x_t - \bar{x})(y_t - \bar{y})}{\sqrt{\sum_{t} (x_t - \bar{x})^2} \sqrt{\sum_{t} (y_t - \bar{y})^2}}
\end{equation}

\textbf{Хөндлөн корреляц ба цагийн хоцрогдол.} Нэг хүний хөдөлгөөн нөгөөгөөсөө тодорхой хугацаагаар хоцорч байгаа эсэхийг тодорхойлоход хөндлөн корреляцыг ашиглана. Хоёр цувааг τ кадраар шилжүүлэн корреляцийг тооцоолж, хамгийн их утга өгсөн \(\tau^{*}\) шилжилтийг цагийн хоцрогдол гэж үзнэ:

\begin{equation}
\rho(\tau) = \frac{\sum_{t} (x_t - \bar{x})(y_{t+\tau} - \bar{y})}{\sqrt{\sum_{t} (x_t - \bar{x})^2} \sqrt{\sum_{t} (y_{t+\tau} - \bar{y})^2}}, \qquad \tau^{*} = \arg\max_{|\tau| \le \tau_{max}} \rho(\tau)
\end{equation}

Энд \(\rho(\tau)\) нь τ кадраар шилжүүлсэн хоёр цувааны Пирсоны корреляц бөгөөд хэвийн болгосон хэлбэртэй тул хөдөлгөөний далайцын ялгаа (жишээлбэл, нэг хүн өвдгөө 12°, нөгөө нь 6° нугалсан) хоцрогдлын тооцоололд нөлөөлөхгүй. Хамгийн их корреляцийн утга \(\rho_{max} = \rho(\tau^{*})\) нь хоёр хөдөлгөөн хэр сайн таарч байгааг харуулдаг тул олдсон хоцрогдолд итгэх эсэхийг шийдэхэд ашиглагдана.

Энэхүү ажилд үений өнцгийн зөрүүгээр тухайн агшны байрлалын ижил төстэй байдлыг, Пирсоны корреляцаар хөдөлгөөний уялдааг, хөндлөн корреляцаар цагийн хоцрогдлыг тус тус үнэлж, эдгээр гурван хэмжүүрийг 3.6-д тайлбарласан байдлаар нэгтгэн синхрончлолын нэгдсэн оноо гаргана.

\section{Технологийн судалгаа}

Систем нь вэб хөтөч (browser) дотор ажиллах вэб аппликейшн болон сонголтоор ашиглах Python HPE серверээс бүрдэх бөгөөд туршилтын үр дүнгийн шинжилгээ болон графикийг Python хэл дээр гаргана. Энэ хэсэгт технологийн хувилбар бүрийг харьцуулж, сонголтын үндэслэл, одоогийн хэрэгжилтийн байдал болон цаашдын хөгжүүлэлтийн чиглэлийг тайлбарлав.

Системийн технологийн бүрэлдэхүүнийг (technology stack) доорх хүснэгтэд нэгтгэв. Сонголт бүрийн үндэслэлийг дараагийн хэсгүүдэд тайлбарлав.

\begin{table}[H]
\centering
\caption{Системийн технологийн бүрэлдэхүүн}
\label{tab:stack}
\small
\begin{tabular}{|>{\raggedright\arraybackslash}p{0.234\linewidth}|>{\raggedright\arraybackslash}p{0.260\linewidth}|>{\raggedright\arraybackslash}p{0.372\linewidth}|}
\hline
\textbf{Давхарга} & \textbf{Технологи} & \textbf{Үүрэг} \\ \hline
Програмчлалын хэл & TypeScript & Вэб аппликейшний үндсэн код \\ \hline
Хөгжүүлэлтийн орчин & Node.js, npm, Vite & Хамаарал удирдах, локал сервер ажиллуулах, бүтээх (build) \\ \hline
HPE загвар & RTMPose-m ба хүн илрүүлэгч (ONNX формат) & Хүн бүрийн гол цэгийг илрүүлэх \\ \hline
Загвар ажиллуулах сан & ONNX Runtime Web & Загварыг вэб хөтөч дотор ажиллуулах \\ \hline
Тооцооллын backend & WebGPU (нөөц: WebAssembly) & Тооцооллыг GPU дээр, боломжгүй үед CPU дээр гүйцэтгэх \\ \hline
Видео ба дүрслэл & HTML5 video, Canvas & Кадр унших, араг яс зурах \\ \hline
Үр дүнгийн формат & MP4, JSON, CSV & Араг яс зурагдсан видео, гол цэг ба синхрончлолын өгөгдөл \\ \hline
Үр дүнгийн шинжилгээ & Python: Pandas, NumPy, Matplotlib & Туршилтын үр дүнг шинжлэх, график гаргах \\ \hline
Вэб хөтөч & Chrome, Edge & WebGPU дэмждэг ажиллах орчин \\ \hline
HPE сервер & Python, FastAPI \cite{fastapi}, Uvicorn & Видеог сервер дээр боловсруулж гол цэгийг буцаах API \\ \hline
Серверийн тооцоолол & ONNX Runtime (CUDA, нөөц: CPU), OpenCV \cite{opencv} & Загварыг NVIDIA GPU дээр ажиллуулах, видеог кадр бүрээр унших \\ \hline
Төслийн бүтэц & npm workspaces (monorepo) & Вэб аппликейшн ба серверийг нэг репозиторид хамтад нь удирдах \\ \hline
\end{tabular}
\end{table}

Загвар сонголтын туршилтад (3.2) нэмэлтээр RTMPose-m, ViTPose-B, YOLO11n загваруудыг ONNX Runtime Web-ээр, MediaPipe Pose-ийг MediaPipe Tasks сангаар ажиллуулсан.

\textbf{Ажиллах орчин.} HPE системийг Python дээр desktop програм эсвэл сервер хэлбэрээр, мөн вэб хөтөч дотор ажиллахаар хэрэгжүүлж болно. Python орчинд хэрэглэгч бүр програм, сан суулгах шаардлагатай бол сервер ашиглавал хэрэглэгчийн видеог сервер рүү илгээх тул нууцлалын асуудал болон серверийн тооцооллын зардал үүснэ. Иймд бүх тооцоолол хэрэглэгчийн төхөөрөмж дээр хийгдэж, видео төхөөрөмжөөс гарахгүй вэб хөтчийн орчныг үндсэн горим болгон сонгов. Харин сул төхөөрөмжтэй хэрэглэгчдэд зориулж видеог сервер дээр боловсруулах нэмэлт горимыг хэрэгжүүлсэн бөгөөд хэрэглэгч аль горимыг ашиглахаа хуудсан дээрээс сонгоно (3.7). Серверийн горимд ч гэсэн сервер зөвхөн гол цэгийг тодорхойлж, синхрончлолын тооцоолол, статистик болон үр дүн хадгалах ажлыг вэб аппликейшн өөрөө гүйцэтгэнэ.

Одоогийн хөгжүүлэлтийн шатанд системийг Node.js орчинд Vite хэрэгслээр локал сервер дээр ажиллуулдаг бөгөөд загваруудын жингийн файлуудыг тусад нь татаж авдаг. Төсөл нь вэб аппликейшн (\texttt{apps/web}) болон Python сервер (\texttt{apps/server}) гэсэн хоёр хэсэгтэй, npm workspaces-ээр удирдагддаг monorepo бүтэцтэй бөгөөд оролтын видео (\texttt{data/}), үр дүн (\texttt{results/}) болон шинжилгээний скриптүүд (\texttt{analysis/}) хоёр хэсэгт нийтлэг байна. Боловсруулалтын үр дүнг хадгалах нь мөн энэ локал сервероор дамждаг. WebGPU технологийн шаардлагаар Chrome эсвэл Edge хөтөч хэрэглэнэ.

\textbf{Тооцооллын backend.} Вэб хөтөч дотор гүн сургалтын загварыг ажиллуулах гурван үндсэн технологи бий.

\begin{table}[H]
\centering
\caption{Вэб хөтөч дэх тооцооллын технологиуд}
\label{tab:backends}
\small
\begin{tabular}{|>{\raggedright\arraybackslash}p{0.148\linewidth}|>{\raggedright\arraybackslash}p{0.123\linewidth}|>{\raggedright\arraybackslash}p{0.295\linewidth}|>{\raggedright\arraybackslash}p{0.254\linewidth}|}
\hline
\textbf{Технологи} & \textbf{Ажиллах төхөөрөмж} & \textbf{Давуу тал} & \textbf{Сул тал} \\ \hline
WebGPU & GPU & Тооцооллын шэйдер (compute shader) дэмждэг тул гүн сургалтын тооцоололд хамгийн хурдан & Шинэ технологи тул бүх хөтөч, төхөөрөмж дэмждэггүй \\ \hline
WebGL & GPU & Бараг бүх хөтөч дэмждэг & График зурахад зориулагдсан тул гүн сургалтын тооцоололд WebGPU-ээс удаан \\ \hline
WebAssembly (WASM) & CPU & Бүх орчинд ажилладаг & Хамгийн удаан \\ \hline
\end{tabular}
\end{table}

Системд ONNX форматын загварууд (YOLO, RTMPose, ViTPose) WebGPU дээр ажиллаж, хөтөч WebGPU дэмжихгүй үед WASM (CPU) руу шилждэг. MediaPipe нь өөрийн GPU горимоор WebGL дээр ажиллаж, боломжгүй үед CPU руу шилждэг. Иймд 3.2-ын хурдны харьцуулалтад MediaPipe бусад загвараас өөр backend дээр хэмжигдсэнийг анхаарах шаардлагатай.

\textbf{Загвар ажиллуулах сан.} Вэб хөтөчид HPE загвар ажиллуулах түгээмэл сангуудыг харьцуулав.

\begin{table}[H]
\centering
\caption{Загвар ажиллуулах сангуудын харьцуулалт}
\label{tab:runtimes}
\small
\begin{tabular}{|>{\raggedright\arraybackslash}p{0.234\linewidth}|>{\raggedright\arraybackslash}p{0.260\linewidth}|>{\raggedright\arraybackslash}p{0.372\linewidth}|}
\hline
\textbf{Сан} & \textbf{Загварын формат} & \textbf{Онцлог} \\ \hline
ONNX Runtime Web & ONNX & PyTorch зэрэг сургалтын орчноос ONNX формат руу хөрвүүлсэн дурын загварыг ажиллуулна; WebGPU, WebGL, WASM backend дэмждэг \\ \hline
TensorFlow.js & TensorFlow & TensorFlow загваруудад тохиромжтой; PyTorch загварыг хөрвүүлэх нь төвөгтэй \\ \hline
MediaPipe Tasks & MediaPipe & Google-ийн бэлэн шийдлүүдийг хялбар ажиллуулна; зөвхөн өөрийн загваруудыг дэмждэг \\ \hline
\end{tabular}
\end{table}

YOLOv8n, YOLO11n, RTMPose-m, ViTPose-B загварууд PyTorch орчинд сургагдсан бөгөөд ONNX формат руу хөрвүүлэх боломжтой тул нэг сангаар бүх загварыг ажиллуулах ONNX Runtime Web-ийг \cite{onnxruntime} сонгов. MediaPipe Pose-ийг харьцуулалтын зорилгоор өөрийн MediaPipe Tasks сангаар ажиллуулсан. TensorFlow.js-ийг харьцуулж үзсэн боловч ашиглаагүй. Загвар сонголтын дараа (3.2) систем RTMPose-m загвар болон түүний хүн илрүүлэгчийг ONNX Runtime Web сангаар ажиллуулна.

\textbf{Видео боловсруулалт.} Бичлэгийг боловсруулахдаа видеог тоглуулахгүйгээр кадр бүрийн тодорхой хугацаа руу шилжиж (seek), тухайн кадр бэлэн болсны дараа загварт оруулдаг. Ингэснээр загвар хэчнээн удаан ажилласан ч нэг ч кадр алгасахгүй. Шилжих хугацааг кадрын эхлэл дээр (\(i/f\)) биш, кадрын голд (эхний кадрын цагийн тэмдэг \(+ (i + 0.5)/f\)) сонгодог. Кадрын яг хил дээр шилжихэд хөтөч заримдаа өмнөх кадрыг харуулдаг бөгөөд үүнээс болж нэг кадр давхардах, нөгөө нь алгасагдах алдаа гардгийг туршилтаар илрүүлж зассан (4.3). Илэрсэн араг ясыг кадр бүр дээр Canvas элементэд зурж, араг яс зурагдсан (annotated) MP4 видео болгон кодчилдог. Үүнээс гадна вэб камераас шууд дүрс авах горим бий бөгөөд энэ горимд кадрыг хөтчийн дэлгэц шинэчлэлтийн давтамжтай уялдуулан боловсруулдаг.

\textbf{Үр дүнгийн шинжилгээ.} Вэб систем гол цэгүүд болон үр дүнг JSON ба CSV файлаар гаргаж, Python скриптүүд тэдгээрийг уншин шинжилгээ хийнэ. Python нь шинжлэх ухааны тооцооллын сангуудаар баялаг бөгөөд Pandas санг өгөгдөл унших, нэгтгэхэд, NumPy санг тоон тооцоололд, Matplotlib санг тайлангийн график зурахад ашиглана. Жишээлбэл, 3.2-ын загваруудын харьцуулалтад Python скрипт загвар бүрийн гол цэгийн JSON файлыг уншиж, консенсусыг тооцоолсон.

\textbf{Хоёр үе шаттай ажиллагаа.} Видеог боловсруулах болон үр дүнг харуулах ажлыг зэрэг гүйцэтгэвэл тооцооллын нөөц их шаардагдах бөгөөд хэрэглэгч бүр хүчирхэг компьютертэй байдаггүй. Сонгосон RTMPose-m загвар нэг кадрт дунджаар 119 мс зарцуулдаг нь секундэд ойролцоогоор 8.4 кадр болно. Кадр руу шилжих, араг яс зурах, MP4 кодчилох зардлыг оруулсан бодит хурд секундэд 5.3 кадр байсан тул 25 кадр/с давтамжтай 1 минутын бичлэгийг боловсруулахад ойролцоогоор 5 минут шаардагдана. Кадрыг нарийн уншдаг болгосны дараах хэмжилтээр RTMPose-m вэб хөтөч дээр (WebGPU) нэг кадрт дунджаар 114 мс, бодит хурд секундэд 5.5 кадр, сервер дээр (CUDA, NVIDIA RTX 2060) 51 мс байв (4.3). Систем нь шүүгч, дасгалжуулагч бичлэгийг дараа нь шинжлэх зориулалттай тул хурднаас илүү нарийвчлалыг чухалчилсан бөгөөд хоёр үе шаттай бүтэц нь боловсруулалт удаан ч үр дүнг тасалдалгүй харуулах боломж олгоно. Иймд систем хоёр үе шаттай ажилладаг.

\begin{enumerate}
\def\labelenumi{\arabic{enumi}.}
\item
  \textbf{Боловсруулах үе шат:} Кадр бүрийг боловсруулж, араг яс зурагдсан MP4 видео болон JSON, CSV үр дүнгийн файлуудыг үүсгэнэ. Энэ үед дэлгэцэнд зөвхөн явцын хувийг харуулна.
\item
  \textbf{Харуулах үе шат:} Бэлэн болсон видеог тоглуулж, тухайн мөчийн синхрончлолын оноо, биеийн хэсэг бүрийн оноо болон цагийн хоцрогдлыг хажууд нь харуулна.
\end{enumerate}

Ингэснээр боловсруулалтын хурд хэрэглэгчийн компьютерээс хамаарч өөр байсан ч харуулах үе шат бүх төхөөрөмж дээр жигд, тасалдалгүй ажиллана.

\textbf{Цаашдын хөгжүүлэлт.} Эцсийн хувилбарт системийг статик вэб хуудас болгон бүтээж (build) интернэтэд байршуулснаар эцсийн хэрэглэгч ямар нэг програм суулгалгүй зөвхөн холбоосоор нэвтэрч ашиглах боломжтой болно. Үүний тулд одоо локал сервероор дамжуулж буй үр дүн хадгалах үйлдлийг хөтчөөс шууд файл татаж авах хэлбэрт шилжүүлнэ.

Мөн гар утасны хөтөч дээр ажиллах боломжийг нэмэлт зорилт болгон авч үзнэ. Гар утасны хөтчүүдийн WebGPU дэмжлэг харилцан адилгүй бөгөөд RTMPose-m загвар хүн бүрийг тусад нь боловсруулдаг тул утсан дээр боловсруулалт мэдэгдэхүйц удаан явагдана. Гэвч системийн гол хэрэглэгчид нь хангалттай хүчин чадалтай компьютер ашиглах шүүгч, дасгалжуулагч бөгөөд хоёр үе шаттай бүтцийн ачаар удаан төхөөрөмж дээр ч үр дүн буруу гарахгүй. Сул төхөөрөмжтэй хэрэглэгчдэд зориулсан бичлэгийг серверт боловсруулах сонголтыг мөн хэрэгжүүлсэн (3.7).
